\documentclass[11pt]{article}
\usepackage[margin=0.9in]{geometry}
\usepackage{microtype}
\usepackage{url}
\usepackage[hidelinks]{hyperref}
\urlstyle{tt}
\setlength{\parskip}{5pt}
\setlength{\parindent}{0pt}
\pagestyle{empty}

\begin{document}

\begin{center}
{\Large\textbf{Resilient Emergency MANETs for Wilderness Safety}}\\[3pt]
{\large One-Page Project Summary}\\[3pt]
{\small Ethan Doherty --- Directed Study (CS 5976) --- Advisors: Prof.~Guevara and Prof.~S.~Basagni --- July 2026}
\end{center}

\vspace{4pt}
\textbf{What this project was.}
Hikers in New Hampshire's mountains regularly lose cell coverage exactly where
emergencies happen. This directed study asked whether a mesh of small, cheap,
solar-powered LoRa radios (Meshtastic) could carry emergency traffic --- SOS
messages, position beacons, short texts --- across that terrain, and set out to
answer it two ways: build and field-test the data-collection system, and
gather ${\geq}2{,}500$ controlled radio measurements to calibrate a
terrain-based coverage model.

\textbf{What happened.}
The first field trial (Mt.~Washington, May) was a systems shakedown: it proved
the collection pipeline and showed that free-space coverage assumptions fail
badly in real terrain --- modeled disagreements up to ${\sim}60$\,dB on
summit-facing links --- which redesigned the field protocol around a
controlled beacon and direct-link-only data. The project's center of gravity
became a purpose-built simulator: two independent implementations (Python and
Rust) of the same statewide mesh model, cross-checked against each other, that
simulate a full year of network life --- radios, batteries, solar, weather,
hikers, emergencies --- in about 3.5 minutes. Its principal finding (a model
result awaiting field calibration): \emph{node survival is governed by how
much time radios spend listening, not by which routing algorithm they run}.
Five different routing algorithms differed by less than 0.3\% in annual
battery depletions, while duty-cycling policies (letting radios sleep between
listen windows) cut depletions 5.5--24$\times$ --- at a quantified cost in
delivery rate and SOS latency, so the result is a design trade, not a free
win. The second field campaign (July) was re-planned in flight after the
beacon radio's hardware failed --- its battery, then its charging port --- so
the calibration dataset was not collected (0 of the 2{,}500 points). The final
field day (Pack Monadnock) still ran the receiver system end-to-end and
measured exactly why public mesh stations cannot substitute for a controlled
beacon: they transmit only 0--4 times per hour, and every packet received
arrived through relays, making signal strength unattributable to any path.
Every change of plan was registered, with cryptographic hashes, before the
data it affected was collected.

\textbf{What it is now.}
A working, documented, reproducible research platform: the two-engine
simulator with a released, hash-locked results set; a terrain-aware
(Longley--Rice) planning layer; a preregistered field protocol with frozen
prediction packs ready to score; a hardened field acquisition rig proven in
the field; and a complete evidence chain --- every number in the final report
traces to a checksummed artifact in the repository
(\url{https://github.com/edoherty2017/resilient-emergency-manets-final}).
One small step remains to close the original calibration goal: a
${\sim}\$25$ replacement beacon board and a single two-radio field day under
the already-frozen protocol. Full detail: \emph{Directed Study Final Report},
2026-07-26.

\end{document}
